\documentclass[../../main.tex]{subfiles}
\begin{document}

\begin{center}
  {\fontsize{48pt}{58pt}\selectfont \textbf{1977}\normalsize\textbf{年度}}
\end{center}

\begin{table}[htb]
  \centering
  \caption{出題テーマと難易度}
  \begin{tabular}{|c|p{6.2cm}|c|c|c|}
    \hline
    \textbf{問題番号} & \textbf{テーマ} & \textbf{思考量} & \textbf{計算量} & \textbf{総合} \\
    \hline
    第1問 & 直交する単位ベクトルの成分を根にもつ方程式 & 易 & 易 & 易 \\
    \hline
    第2問 & 三次関数と座標軸に囲まれる面積の最大化 & 易 & 易 & 易 \\
    \hline
    第3問 & 三角関数の高階導関数と交点での接線 & 易 & 易 & 易 \\
    \hline
    第4問 & サイコロの振り直しの最適戦略 & 易 & 易 & 易 \\
    \hline
    第5問 & 周長を固定した三角形・四辺形の最大面積 & 普 & 普 & 普 \\
    \hline
    第6問 & 微分不等式から得る関数の零点 & 難 & 難 & 難 \\
    \hline
  \end{tabular}
\end{table}

\begin{table}[htb]
  \centering
  \caption{解答}
  \renewcommand{\arraystretch}{1.4}
  \begin{tabular}{|c|c|p{9cm}|}
    \hline
    \textbf{問題番号} & \textbf{小問} & \textbf{答え} \\
    \hline
    第1問 & -- & $x^2\pm\dfrac{\sqrt{15}}{3}x+\dfrac{1}{3}=0$ \\
    \hline
    第2問 & -- & $a=1$ \\
    \hline
    第3問 & -- & 交点での2接線は直交する（証明問題） \\
    \hline
    第4問 & (1) & 1回目が$3$以下なら振り直す \\
    \cline{2-3}
    & (2) & 1回目が$4$以下なら振り直し，2回目が$3$以下ならさらに振り直す \\
    \hline
    第5問 & (1) & 面積最大は二等辺三角形（証明問題） \\
    \cline{2-3}
    & (2) & 面積最大は正方形（証明問題） \\
    \hline
    第6問 & (i) & $g(a)>0$ \\
    \cline{2-3}
    & (ii) & $a\le y\le a+\dfrac{f(a)}{g(a)}$に$f(y)=0$を満たす$y$が存在する \\
    \hline
  \end{tabular}
\end{table}

\paragraph{試験形態}
6問．試験時間は要確認．

\end{document}
